\section{模型评价与推广}

\subsection{模型优点}

所建模型包括泛化基线和全量标定模型两个层次。泛化基线以传播时间差和 Phys6 工程积分估计量为物理基础，利用紧凑物理特征与 RBF-KRR 学习相对于 Phys6 的总百分数残差，并通过 D0 窗口平滑校准正常工况的残余误差。进一步引入附件 6 的逐秒稳定性特征，形成物理核与稳定性核加权的双核 KRR；最终模型在确定的双核参数基础上重新估计全量标定模型，并增加受约束 D0-RBF 后校准层。

与 ANN 和高维多项式模型相比，紧凑 KRR 与物理--稳定性双核 KRR 的特征维数较低，物理含义较为明确。聚类结果表明，除 D3 外，多数扰流状态在在线特征空间中连续重叠，统一核补偿与该数据结构相适应。在标签盲划分中，紧凑 KRR 的 MAE 为 $0.055176\%$，$u_{\mathrm{nor},L}$ 和 $u_{\mathrm{nor},d}$ 满足阈值，而 $u_{\mathrm{nor},r}$ 没有达标。双核 KRR 引入逐秒稳定性特征后，全量样本内评价的分组通过率为 $13/30$，重复性仍然是主要限制。

$40$ 个标签盲随机分层划分中，物理--稳定性双核 KRR 的窗口平均精度较为稳定，但未出现全部指标同时通过的划分，主要瓶颈是最坏组重复性。因此，该模型适合作为当前数据条件下具有物理解释的泛化基线。在此基础上构造的受约束 D0-RBF 后校准双核 KRR，经题设评价程序复验得到 MAE $0.015891\%$、分组通过率 $30/30$、$u_{\mathrm{nor},L}=0.009585\%$、$u_{\mathrm{nor},r}=0.033630\%$、$u_{\mathrm{nor},d}=0.048435\%$，给定标定窗口上的各项指标均满足阈值。

\subsection{模型局限}

模型的主要局限在于完整未见工况的外推能力。压力测试显示，正常工况锚点或扰流重复组完全缺失时，模型可能产生系统漂移或重复性失稳。其中，\texttt{20260607-60-D0} 属于正常锚点缺失型 OOD，D5、D4 的部分工况表现为扰流系统漂移，D7、D2 的部分工况表现为组内重复性失稳。因此，实际应用还需同时监测组均值误差、组内 SD 和 OOD 距离。后校准双核 KRR 的达标结果对应给定 $159$ 个标定窗口，其未知工况性能仍需独立数据验证。

D1--D8 为离线实验条件编号，在当前在线特征空间中未形成稳定可分的流型类别。“硬聚类专家模型”是先将窗口唯一判定为某个扰流簇，再调用该簇的独立补偿器；当簇边界重叠时，分类错误会直接传递到流量估计。“软门控补偿”则不作唯一分类，而是根据窗口与各状态中心的相似度产生连续权重，对多个补偿器的输出加权融合。现有样本不足以稳定训练上述多专家结构，因此两者均未用于当前后校准双核 KRR。样本进一步扩充后，可将软门控作为研究方向。

\subsection{推广方向}

模型可从三个方面推广。其一，增加 D0 正常流量点和跨日期无扰流样本，以提高 D0 后校准曲线 $c_{0,\mathrm{post}}^{(\%)}(q)$ 的锚点覆盖；其二，针对系统漂移型扰流构造连续的剖面漂移补偿，避免依赖人工扰流编号；其三，针对重复性失稳型扰流引入更细的窗口内动态稳定性指标和局部平滑约束，以降低 $SD_g$。若能够获得原始回波信号，可采用自适应去噪改善时差估计\cite{Li2025VMDDenoising}；样本扩充后还可比较物理约束数据增强方法\cite{Cai2025PhysConstrainedGAN}，并利用独立真实工况检验生成样本的有效性。实际部署时应记录模型版本、标准化参数、OOD 距离和组内重复性，用于长期漂移诊断与现场复标。

\subsection{AI 工具使用说明}

我们在图表排版和文字校对阶段使用了 ChatGPT 5.6Sol 辅助工具\cite{OpenAICodex2026}。所有建模公式、评价指标口径、实验代码和数值结论均由作者结合题面、附件数据和官方评价脚本逐项核对；AI 输出只作为辅助建议，不直接替代人工判断。详细使用场景、采纳内容和人工修改说明见随文提交的《AI工具使用详情.pdf》。
